% Problem-section file following ../_template.tex.
\section{Constant trace-distance separability testing}
\label{sec:problem-36}

\paragraph{Problem.}
What is the computational complexity of testing bipartite separability with
a constant trace-distance promise gap?  For local dimension $d$, define
\begin{equation}
  \operatorname{Sep}(d,d)
  :=\operatorname{conv}\!\left\{
    \alpha\otimes\beta:
    \alpha,\beta\in\mathcal D(\mathbb C^d)
  \right\}.
  \label{eq:p36-separable-set}
\end{equation}
Fix a constant $0<\varepsilon_0<1$.  Given a rational description of
$\rho\in\mathcal D(\mathbb C^d\otimes\mathbb C^d)$, promised that exactly
one of the following alternatives holds,
\begin{equation}
  \begin{aligned}
    \text{YES:}\quad&\rho\in\operatorname{Sep}(d,d),\\
    \text{NO:}\quad&
      \inf_{\sigma\in\operatorname{Sep}(d,d)}
      \frac12\lVert\rho-\sigma\rVert_1\geq\varepsilon_0,
  \end{aligned}
  \label{eq:p36-trace-promise}
\end{equation}
decide which case in Eq.~\eqref{eq:p36-trace-promise} holds.  Is there an
algorithm polynomial or quasipolynomial in $d$?  More generally, when the gap
$\varepsilon$ is part of the input, determine the optimal dependence of the
complexity on $d$ and $\varepsilon$.

\subsection*{Status}

Unsolved

\subsection*{Source}

Harrow and Montanaro explicitly pose constant-gap weak membership for
separability in trace norm as an open complexity problem
\sourcecite{ref:p36-harrow}{HM13}.

\subsection*{Progress}

\begin{itemize}
  \item Harrow and Montanaro formulate the constant trace-distance promise in
  Eq.~\eqref{eq:p36-trace-promise} explicitly and relate it to estimating
  acceptance probabilities in $\mathsf{QMA}(2)$
  \sourcecite{ref:p36-harrow}{HM13}.

  \item Weak membership for the set in
  Eq.~\eqref{eq:p36-separable-set} is strongly $\mathsf{NP}$-hard when the
  promised distance is inverse-polynomial in the input dimension
  \sourcecite{ref:p36-gharibian}{Gha10}.  This hardness regime does not
  classify the fixed constant gap in Eq.~\eqref{eq:p36-trace-promise}.

  \item A symmetric-extension algorithm has running time
  \begin{equation}
    \exp\!\left(
      O\!\left(\varepsilon^{-2}(\log d)^2\right)
    \right)
    \label{eq:p36-extension-runtime}
  \end{equation}
  for Euclidean distance or an operational $\mathsf{LOCC}$ norm
  \sourcecite{ref:p36-brandao}{BCY11}.  The bound in
  Eq.~\eqref{eq:p36-extension-runtime} does not hold as stated for trace
  distance.

  \item For each fixed constant gap, randomized polynomial time is now known
  for Euclidean-norm weak membership
  \sourcecite{ref:p36-malavolta}{Mal26}.  Constant trace distance can coexist
  with Euclidean distance that vanishes with $d$, so this result does not
  solve Eq.~\eqref{eq:p36-trace-promise}.
\end{itemize}

\subsection*{References}

\begin{enumerate}[label={},leftmargin=4.8em,labelsep=0.5em]
  \footnotesize
  \item[\textup{[HM13]}]\label{ref:p36-harrow}
  A. W. Harrow and A. Montanaro, ``Testing Product States, Quantum
  Merlin--Arthur Games and Tensor Optimization,''
  \emph{Journal of the ACM} \textbf{60}, Article 3 (2013).
  \href{https://doi.org/10.1145/2432622.2432625}%
       {doi:10.1145/2432622.2432625};
  \href{https://arxiv.org/abs/1001.0017}{arXiv:1001.0017}.

  \item[\textup{[Gha10]}]\label{ref:p36-gharibian}
  S. Gharibian, ``Strong NP-Hardness of the Quantum Separability Problem,''
  \emph{Quantum Information and Computation} \textbf{10}, 343--360 (2010).
  \href{https://arxiv.org/abs/0810.4507}{arXiv:0810.4507}.

  \item[\textup{[BCY11]}]\label{ref:p36-brandao}
  F. G. S. L. Brand\~ao, M. Christandl, and J. Yard,
  ``A Quasipolynomial-Time Algorithm for the Quantum Separability Problem,''
  in \emph{Proceedings of the 43rd Annual ACM Symposium on Theory of
  Computing}, 343--351 (2011).
  \href{https://doi.org/10.1145/1993636.1993683}%
       {doi:10.1145/1993636.1993683};
  \href{https://arxiv.org/abs/1011.2751}{arXiv:1011.2751}.

  \item[\textup{[Mal26]}]\label{ref:p36-malavolta}
  G. Malavolta, ``Quantum Separability in Polynomial Time,''
  arXiv:2607.23773 (2026).
  \href{https://arxiv.org/abs/2607.23773}{arXiv:2607.23773}.
\end{enumerate}

\subsection*{Comment}

The constant-gap trace-norm task in
Eq.~\eqref{eq:p36-trace-promise} is posed explicitly in Section~4.2,
item~14 of \sourcecite{ref:p36-harrow}{HM13}.  The norm is essential: the
constant-gap Euclidean problem has been solved, but that result does not
classify trace-norm weak membership.

\subsection*{Tag}

Quantum separability; Quantum complexity theory; Convex optimization;
Semidefinite programming; Entanglement theory

\subsection*{ID}

\texttt{op\_ec184fc49232a9c0}
